A primeira metade de 2025 foi um período extensivamente determinante para os sentidos que o correr da vida tomou;
    sobrepunham-se eventos impactantes, contínuos e incisivos em minha caminhada, desviantes de mim, num tempo que me atropelava e seguia em frente como um trem que eu sigo correndo para alcançar:
        foi nesse bosque, em meio a uma vasta e mestiça colheita, que me encontrei com deus e o diabo.

Ambos a provocar insensibilidades aos dizeres do bosque que eu buscava -- e sigo buscando -- transformar em novos modos de escutar o mundo;
    se o kit-de-odores é um artifício para cultivar narizes em corpos desarticulados a sutis contrastes olfativos, deus e o diabo foram então cúmplices no cultivo de órgãos sensíveis à continuidade espiralar entre o sagrado e o profano:
        ora os via separados por longas distâncias, ocupando territórios em margens opostas dentro no bosque -- ora os encontrava habitando simultaneamente o espaço inteiro e os seus viventes todos, inclusive a mim.

Como uma moeda;
    em lados opostos, "cara" e "coroa" nunca se encontram, jamais se olham nos olhos -- mas conformam existência, significado e agência tanto um ao outro, quanto ao pedaço de metal lhes comporta reciprocamente a vida:
        um verdadeiro Totem.

---

Em nosso primeiro encontro, deus me pôs a pensar na morte, a imaginar o ano exato em que retornarei ao pó e as palavras que desejaria deixar em minha lápide, como um pai rigoroso que não deixa se esquecer da dureza e finitude da vida;
    já o diabo me apresentou diversas trilhas, categorizou e descreveu os caminhos e destinos de cada uma, compartilhou conselhos e segredos como um velho sábio a acolher um bando de jovens se aventurando pelo bosque:
        diferente do que se poderia esperar, no entanto, era deus que se disfarçava e escondia seus truques, e o diabo é que se apresentava com cintilante transparência.

Sagrado é o caminho, dizia o todo-poderoso, é a linha costurada por cada um desde o próprio nascimento até o instante de sua morte;
    profano é o que está fora das linhas consagradas, alertava o demônio, insistindo nos perigos e na imoralidade de se abrirem e percorrerem novas trilhas pelo bosque:
        uma leitora moderna e sagaz logo apontaria a redundância dos papéis, afinal, basta um norte para que a bússola nos aponte a direção -- mas aquelas versadas no politeísmo têm a disposição uma miríade de instrumentos de navegação, e sabem que sem o sul ficam presas sempre aos mesmos caminhos do norte.

Defronte ao mesmo mato fechado no meio do meu caminho, deus e o diabo indicavam sentidos opostos;
    olhe para dentro -- não olhe para fora:
        o mesmo mato que é fora é dentro, que é sagrado é profano.

